\chapter{Female Reproductive System}

\label{ch:ch52_female_reproductive}

\section{What imaging covers}
Female reproductive imaging evaluates the uterus, fallopian tubes, ovaries, and surrounding pelvic structures. The most common investigations are pelvic ultrasound (transabdominal and transvaginal), pelvic MRI, and CT for acute or oncologic indications.

\section{When is imaging ordered?}
\begin{itemize}
    \item \textbf{Pelvic pain:} Ectopic pregnancy, ovarian torsion, endometriosis, pelvic inflammatory disease.
    \item \textbf{Abnormal bleeding:} Fibroids, adenomyosis, polyps, endometrial hyperplasia, or malignancy.
    \item \textbf{Palpable mass:} Ovarian cyst, fibroid, or ovarian tumor.
    \item \textbf{Infertility:} Ovarian reserve assessment, uterine anomalies, tubal patency.
    \item \textbf{Known cancer:} Staging of endometrial, cervical, or ovarian cancer.
    \item \textbf{Pregnancy complications:} Viability, location, bleeding.
\end{itemize}

\section{Transabdominal vs transvaginal ultrasound}

\begin{table}[h]
\centering
\caption{Comparison of ultrasound approaches}
\label{tab:ta_vs_tv}
\begin{tabular}{@{}p{0.18\textwidth}p{0.31\textwidth}p{0.37\textwidth}@{}}
\hline
\textbf{Feature} & \textbf{Transabdominal} & \textbf{Transvaginal} \\
\hline
Best for & Large masses, pregnancy, pregnancy & Early pregnancy, ovarian detail, endometrium \\
Bladder & Must be full & Not required (may be empty) \\
Image resolution & Lower & Higher (probe closer to organs) \\
Operator skill & Basic & Requires specific training \\
Patient comfort & Non-invasive & Mildly uncomfortable; consent required \\
Probe placement & On abdomen & Inside vagina \\
\hline
\end{tabular}
\end{table}

\begin{pearlbox}[Transvaginal ultrasound and early pregnancy]
Transvaginal ultrasound detects an intrauterine pregnancy 1--2 weeks earlier than transabdominal scanning. A gestational sac should be visible at $\geq$5 weeks, and a yolk sac at $\geq$5.5 weeks. If no intrauterine pregnancy is seen with a $\beta$-hCG above the discriminatory threshold (typically 1500--2000~mIU/mL), ectopic pregnancy must be actively excluded.
\end{pearlbox}

\section{Ovarian masses: the O-RADS system}
The Ovarian-Adnexal Reporting and Data System (O-RADS) classifies ovarian masses by risk of malignancy and guides management:

\begin{table}[h]
\centering
\caption{O-RADS classification}
\label{tab:orads}
\begin{tabular}{@{}p{0.12\textwidth}p{0.31\textwidth}p{0.14\textwidth}p{0.28\textwidth}@{}}
\hline
\textbf{Score} & \textbf{Description} & \textbf{Malignancy risk} & \textbf{Management} \\
\hline
0 & Incomplete assessment & --- & Repeat scan or MRI \\
1 & Functional cyst $<$3~cm & $<$1\% & No follow-up needed \\
2 & Simple unilocular anechoic, thin wall, no solid & $<$1\% & Follow-up or surgery based on size \\
3 & Multilocular, no solid component, thin septa & $<$3\% & Surgery if $>$10~cm; follow-up if smaller \\
4 & Solid component or thick septa & 3--50\% & Surgery + oncology referral \\
5 & Clearly malignant features & $>$50\% & Oncology referral \\
\hline
\end{tabular}
\end{table}

\section{Common ovarian cysts}

\begin{table}[h]
\centering
\caption{Types of ovarian cysts}
\label{tab:ovarian_cysts}
\begin{tabular}{@{}p{0.22\textwidth}p{0.35\textwidth}p{0.29\textwidth}@{}}
\hline
\textbf{Cyst type} & \textbf{US features} & \textbf{Clinical note} \\
\hline
Corpus luteum & Thick-walled, may contain echogenic material & Can cause positive $\beta$hCG if very large \\
Dermoid (mature cystic teratoma) & Complex: echogenic foci (fat), calcification, ``tip of iceberg'' & Most common ovarian germ cell tumor \\
Endometrioma & Homogeneous low-level echoes (``ground glass'') & Bilateral in 40\%; associated with endometriosis \\
Hemorrhagic cyst & Reticular internal echoes (``fishnet'' pattern) & Resolves in 6--8 weeks; do not biopsy \\
Simple follicular & Anechoic, thin wall, posterior enhancement & Functional; resolves in 1--3 cycles \\
\hline
\end{tabular}
\end{table}

\begin{pearlbox}[Never biopsy a dermoid]
Mature cystic teratomas (dermoid cysts) contain fat, hair, teeth, and other ectodermal elements. They should not be aspirated percutaneously due to risk of chemical peritonitis from spillage of sebaceous material. Surgical excision is the management of choice.
\end{pearlbox}

\section{Uterine fibroids (leiomyomas)}
Fibroids are the most common uterine tumors, found in up to 70\% of women over 50. Imaging features:
\begin{itemize}
    \item \textbf{Ultrasound:} Well-circumscribed, heterogeneous, hypoechoic mass within the myometrium. Calcification produces acoustic shadowing.
    \item \textbf{MRI:} T2-hypointense mass (fibroid) with or without degeneration. Useful for mapping fibroids pre-surgery or pre-uterine artery embolization (UAE).
    \item \textbf{Classification by location:} Submucosal (into cavity---causes bleeding), intramural (within wall), subserosal (projecting outward).
\end{itemize}

\section{Adenomyosis}
Adenomyosis is endometrial tissue within the myometrium, causing painful, heavy periods.

\begin{table}[h]
\centering
\caption{Adenomyosis vs fibroid}
\label{tab:adeno_vs_fibroid}
\begin{tabular}{@{}p{0.20\textwidth}p{0.42\textwidth}p{0.30\textwidth}@{}}
\hline
\textbf{Feature} & \textbf{Adenomyosis} & \textbf{Fibroid} \\
\textbf{Junctional zone} & $>$12~mm (pathological) & Normal unless distorted \\
\textbf{Margins} & Ill-defined & Well-defined \\
\textbf{Shape} & Diffuse/ovoid & Round/lobulated \\
\textbf{T2 MRI} & High-signal foci within thickened junctional zone & Low-signal mass \\
\textbf{Vascularity} & Peripheral vascularity & Peripheral ``ring of fire'' \\
\hline
\end{tabular}
\end{table}

\section{Ectopic pregnancy}
Ectopic pregnancy (implantation outside the uterus, usually in the fallopian tube) is a life-threatening emergency. Imaging findings:

\begin{itemize}
    \item \textbf{Ultrasound:} No intrauterine pregnancy with elevated $\beta$-hCG above discriminatory level. Adnexal mass with or without cardiac activity. ``Tubal ring sign''---a hyperechoic ring surrounding a gestational sac. Free pelvic fluid (suggests rupture).
    \item \textbf{CT:} Used when ruptured ectopic is suspected and the patient is haemodynamically unstable. Shows complex adnexal mass, haemoperitoneum, and absent intrauterine pregnancy.
\end{itemize}

\begin{pearlbox}[Discriminatory level]
When $\beta$-hCG reaches 1500--2000~mIU/mL, an intrauterine pregnancy should be visible on transvaginal ultrasound. If it is not, the differential is ectopic pregnancy, recent miscarriage, or early viable intrauterine pregnancy. Do not rely on a single scan---correlate with serial $\beta$-hCG and repeat imaging.
\end{pearlbox}

\section{Ovarian torsion}
Ovarian torsion is twisting of the ovary on its vascular pedicle, compromising blood flow.
\begin{itemize}
    \item \textbf{Risk factors:} Ovarian cyst or mass $>$4~cm, ovarian hyperstimulation, pregnancy.
    \item \textbf{Ultrasound:} Enlarged ovary with peripheral follicles, absent or reduced Doppler flow (venous flow lost first, then arterial), twisted pedicle (``whirlpool sign'' on color Doppler), and free fluid.
    \item \textbf{Important:} Normal Doppler flow does \emph{not} exclude torsion---the ovary has dual blood supply, and partial torsion may show preserved flow.
\end{itemize}

\section{Endometrial thickness}
Endometrial thickness is measured on the sagittal transvaginal view in the mid-sagittal plane. It is a key marker for endometrial pathology:

\begin{table}[h]
\centering
\caption{Endometrial thickness thresholds}
\label{tab:endometrium}
\begin{tabular}{@{}p{0.30\textwidth}p{0.13\textwidth}p{0.44\textwidth}@{}}
\hline
\textbf{Clinical scenario} & \textbf{Threshold} & \textbf{Significance} \\
\hline
Postmenopausal bleeding & $>$4~mm & Requires further investigation (biopsy/hysteroscopy) \\
Postmenopausal on HRT & $>$12~mm & Requires investigation \\
Premenopausal (proliferative) & 4--8~mm & Normal \\
Premenopausal (secretory) & 8--16~mm & Normal \\
\hline
\end{tabular}
\end{table}

\section{MRI for pelvic masses: when to recommend}
MRI is recommended when:
\begin{itemize}
    \item Ultrasound findings are indeterminate (solid vs cystic, origin uncertain)
    \item The mass appears to arise from the uterus but the extent is unclear
    \item Pre-surgical planning for fibroid mapping or ovarian cancer staging
    \item Differentiating adenomyosis from fibroids
    \item Assessing deep infiltrating endometriosis
    \item Characterising complex adnexal masses (e.g., distinguishing endometrioma from hemorrhagic cyst)
\end{itemize}

\section{Reporting checklist}
A structured female reproductive report should include:
\begin{enumerate}
    \item Uterus: size, position (ante- vs retroverted), endometrial thickness and echogenicity
    \item Myometrium: homogeneous vs fibroids (size, number, location) vs adenomyosis
    \item Cervix: length, endocervical canal
    \item Ovaries: size, follicular pattern, dominant cysts, solid masses (O-RADS score)
    \item Adnexa: masses, free fluid, tubal dilatation
    \item Pregnancy status (if relevant): location, viability, gestational age
    \item Pelvic free fluid
    \item Comparison with prior imaging
\end{enumerate}

\clearpage
